\begin{example}[The zero morphism]
For $0:\mathcal F\to\mathcal G$, $\ker 0=\mathcal F$ and $\im 0=0$. The image presheaf is already the zero sheaf, so no repair is needed.
\end{example}

\begin{example}[The identity morphism]
For $\id_{\mathcal F}$, the kernel is $0$, the image is $\mathcal F$, and the cokernel is $0$. This is the sheaf version of the most elementary group calculation.
\end{example}

\begin{example}[Multiplication by an integer]
On the constant sheaf $\underline{\mathbb Z}$, multiplication by $n$ has kernel $0$ for $n\neq0$ and cokernel $\underline{\mathbb Z/n\mathbb Z}$. Stalkwise this is the ordinary map $\mathbb Z\xrightarrow{n}\mathbb Z$.
\end{example}

\begin{example}[Restriction to an open subset]
If $j:U\hookrightarrow X$ and $\varphi:\mathcal F\to\mathcal G$, then $(\ker\varphi)|_U=\ker(\varphi|_U)$. The same holds for image and cokernel because these are characterized locally.
\end{example}

\begin{example}[A skyscraper kernel]
Let $i_x(A)$ be the skyscraper sheaf at a closed point $x$ and let $A\to B$ be a group map. The induced skyscraper morphism has kernel $i_x(\ker(A\to B))$, image $i_x(\im(A\to B))$, and cokernel $i_x(\coker(A\to B))$.
\end{example}

\begin{example}[Quotient of skyscraper sheaves]
If $B\subseteq A$, then $i_x(A)/i_x(B)\cong i_x(A/B)$. Every nonzero stalk occurs only at $x$, so the quotient calculation is completely pointwise.
\end{example}

\begin{example}[Discrete spaces]
If $X$ is discrete, a sheaf is a product of its stalks: $\mathcal F(U)\cong\prod_{x\in U}\mathcal F_x$. Kernels, images, and quotients are then computed componentwise and the naive sectionwise constructions are already sheaves.
\end{example}

\begin{example}[A kernel as vanishing locus of a restriction]
Let $V\subseteq U$ be open. For the restriction morphism of sheaves $\mathcal C^0_U\to j_*\mathcal C^0_V$, the kernel consists of continuous functions whose restriction to $V$ is zero. This condition is local and therefore automatically defines a sheaf.
\end{example}

\begin{example}[A quotient with no correction needed]
On a discrete two-point space, let $\mathcal F$ have stalks $\mathbb Z$ and let $2\mathcal F\subseteq\mathcal F$. Then $(\mathcal F/2\mathcal F)(U)$ is exactly the product of copies of $\mathbb Z/2\mathbb Z$ over points of $U$.
\end{example}

\begin{example}[Locally constant functions modulo $n$]
For the locally constant sheaf $\underline{\mathbb Z}$, the quotient by $n\underline{\mathbb Z}$ is $\underline{\mathbb Z/n\mathbb Z}$. This can be checked on stalks or by locally constant representatives.
\end{example}

\begin{example}[Image of an inclusion]
If $\mathcal F'\hookrightarrow\mathcal F$ is literally a subsheaf inclusion, then the image sheaf is $\mathcal F'$. Its sectionwise image is already a sheaf because the map identifies $\mathcal F'(U)$ with a subgroup of $\mathcal F(U)$ compatible with gluing.
\end{example}

\begin{example}[Projection from a direct sum]
For $\mathcal F\oplus\mathcal G\to\mathcal G$, the kernel is $\mathcal F$, the image is $\mathcal G$, and the quotient by the kernel is canonically $\mathcal G$.
\end{example}

\begin{example}[Diagonal and difference]
The diagonal map $\Delta:\mathcal F\to\mathcal F\oplus\mathcal F$ has cokernel canonically isomorphic to $\mathcal F$ via $(a,b)\mapsto b-a$. This can be verified sectionwise because the displayed difference gives a global splitting.
\end{example}

\begin{example}[Functions vanishing at a point]
Let $X$ be a space and $x\in X$. The morphism $\mathcal C^0_X\to i_x(\mathbb R)$ taking a function to its value germ at $x$ has kernel the sheaf of continuous functions vanishing at $x$ in the stalk sense. The quotient has stalk $\mathbb R$ at $x$ and $0$ elsewhere.
\end{example}

\begin{example}[A locally liftable section]
Suppose $g\in\mathcal G(U)$ and for each $x\in U$ there is $V_x$ and $f_x\in\mathcal F(V_x)$ with $\varphi(f_x)=g|_{V_x}$. Then $g$ belongs to the image sheaf even if the $f_x$ do not agree on overlaps and no global $f$ exists.
\end{example}

\begin{example}[Quotient classes are local]
A section of $\mathcal F/\mathcal F'$ over $U$ may be represented by quotient classes $[s_i]$ on an open cover whose differences $s_i-s_j$ lie in $\mathcal F'(U_i\cap U_j)$. The representatives need not glue in $\mathcal F$; their quotient classes do.
\end{example}

\begin{example}[Cokernel of an inclusion]
For $i:\mathcal F'\hookrightarrow\mathcal F$, $\coker i$ is exactly $\mathcal F/\mathcal F'$. The sectionwise quotient may fail to be a sheaf, which is why the definition uses sheafification.
\end{example}

\begin{example}[Kernel and stalks]
If a section $s\in\mathcal F(U)$ has germ $s_x\in\ker\varphi_x$ for every $x\in U$, then $\varphi(s)$ has zero germ everywhere. Since $\mathcal G$ is a sheaf, $\varphi(s)=0$, so $s\in(\ker\varphi)(U)$.
\end{example}

\begin{example}[Image and stalks]
If every germ $g_x$ of $g\in\mathcal G(U)$ lies in $\im\varphi_x$, then for each $x$ choose a local representative lift. Hence $g$ is locally in the image and therefore lies in $(\im\varphi)(U)$.
\end{example}

\begin{example}[First isomorphism theorem]
For any $\varphi:\mathcal F\to\mathcal G$, the canonical map $\mathcal F/\ker\varphi\to\im\varphi$ is an isomorphism. On stalks it is the ordinary first isomorphism theorem for abelian groups.
\end{example}
